\begin{figure}
    \centering
    % \includegraphics[width=\textwidth,height=0.35\textwidth]{example-image-duck}
    \begin{overpic}[width=0.64\textwidth]{figures/extrapolation/extrapolation_illust}
    % \put(4, 49){\footnotesize\xgray{(a)}}
    \end{overpic}
    \hfill
    \begin{overpic}[height=0.34\textwidth,trim={1cm 0.15cm 0 0}]{figures/extrapolation/extrapolation}
    \end{overpic}
    \venueifelse{\vskip -0.5\baselineskip}{\vskip -0.25\baselineskip}
    \caption{\textbf{Selective underfitting.} (a) Training distribution concentrates in a \coloredul{xblue}{small region} of data space, where the model learns the \xxgreen{empirical score function $\rvs_\star$}. At inference, \coloredul{xred}{sampling trajectories} go beyond this region, where predictions are not directly \xxgreen{\emph{\textbf{supervised}}} and must be \xsienna{\emph{\textbf{extrapolated}}}. (b) The plotted distances reflect how far \coloredul{xred}{sampling trajectories} are from the \coloredul{xblue}{supervision region}.}
    \label{fig:extrapolation}
    \phantomsubcaption\label{fig:extrapolation-a}
    \phantomsubcaption\label{fig:extrapolation-b}
    % \neuripsonly{\vskip -0.5\baselineskip}
    \vskip -\baselineskip
\end{figure}